\question[4]
Betrachte die folgende Tabelle:

\begin{center}
\begin{tabular}{|c|l|l|}
\hline
KundenID & Name & Wohnort \\
\hline
101 & Anna Meier & Stuttgart \\
102 & Max Schulz & Böblingen \\
103 & Lea Weber & Sindelfingen \\
\hline
\end{tabular}
\end{center}

Erläutere die Begriffe \textbf{Schlüssel} und \textbf{Primärschlüssel}.
Bestimme außerdem den Primärschlüssel der Tabelle und begründe deine Wahl.

\begin{solutionbox}{4cm}
Ein Schlüssel ist ein Attribut oder eine Attributkombination, mit der ein Datensatz eindeutig identifiziert werden kann.

Ein Primärschlüssel ist der ausgewählte Schlüssel, der jeden Datensatz einer Tabelle eindeutig kennzeichnet.

Der Primärschlüssel der Tabelle ist \textbf{KundenID}, da jede KundenID nur einmal vorkommt und jeden Kunden eindeutig identifiziert.

Die Attribute Name und Wohnort sind ungeeignet, da mehrere Kunden denselben Namen oder Wohnort besitzen könnten.
\end{solutionbox}